\section{One original joint-harmonic map for every parameter}
\label{sec:joint-harmonic-io}

This chapter reconstructs the simultaneous producer
\texttt{GlobalTheory.JointHarm.exists\_ioP} at source commit
\texttt{7cbf7a0c32fca1badf05931879a37c3f8a2b03b4}.
The formal source is taken as correct; no Lean replay is asserted.
Every assertion below is a human proof on the stated original inputs.
In particular, constructing the map does not construct the original
\(R\), its coefficient tie, its discrete-series models or its Weil data.

We keep the source's conventions.  Inner products are conjugate-linear
in the first argument.  For a sign function
\(\epsilon:\{0,\ldots,m-1\}\to\{+,-\}\), put
\[
 J_\epsilon=\operatorname{diag}(\epsilon_a),\qquad
 K_\epsilon=\{k:k^*k=1,\ kJ_\epsilon=J_\epsilon k\},\qquad
 K_\epsilon^{\mathbb C}=\{g\in\operatorname{GL}_m(\mathbb C):
                         gJ_\epsilon=J_\epsilon g\}.
 \tag{JIO.1}
\]
Thus the last two groups are the compact and complex block groups,
including empty blocks.  A holomorphic representation in this chapter
means exactly the source's \texttt{IsHolRep}: a sesquilinear form \(B\)
with \(B(x,x)=0\Rightarrow x=0\), the representation identities on
\(K_\epsilon^{\mathbb C}\), the adjoint identity
\(B(\pi(g)x,y)=B(x,\pi(g^*)y)\), and holomorphy of every scalar
coefficient \(B(x,\pi(g)y)\) there.  A bottom vector of weight \(\mu\)
is fixed by every lower-unitriangular matrix in this block group and
has diagonal unitary character
\(\exp(i\sum_a\mu_a\theta_a)\).
The words \texttt{IsHolFam} and \texttt{IsBotFam} mean the corresponding
product-group identities and one-place conditions.

Write \(I=D.\mathrm{what}_2.\mathrm{Tau}\).  All products over \(I\)
are finite.  In the polynomial construction \(m=W.\mathrm{rank}\),
\(n=V.\mathrm{rank}\), and
\(\mathcal P=\mathbb C[X_{\tau,a,b}:\tau\in I,a<m,b<n]\)
is the literal \texttt{PolF D m n}.  Later \(V\) and \(W\) are the
stored \(\mathrm{Vc}(D,i)\), \(\mathrm{Wc}(D,i)\), with their original
frames.  No change of an underlying polynomial or coefficient space
is implicit.

% BEGIN JIO1
\paragraph{JIO1: linear substitution and the exact factorial adjoint.}
\label{claim:JIO1}
For any finite variable set \(E\), define
\[
 B_F(P,Q)=\sum_{\alpha}\overline{P_\alpha}Q_\alpha\alpha!,
 \quad \alpha!=\prod_{e\in E}\alpha_e!,
 \qquad L(A)X_e=\sum_f A_{fe}X_f .
 \tag{JIO.2}
\]
Then \(B_F\) is positive definite,
\(L(AA')=L(A)L(A')\), and
\[
                  B_F(L(A)P,Q)=B_F(P,L(A^*)Q).
 \tag{JIO.3}
\]
These statements hold on all polynomials, with no completion.

\emph{Proof.}
The sum in \(B_F(P,P)\) consists of the nonnegative real terms
\(|P_\alpha|^2\alpha!\), and every \(\alpha!\) is positive.  Thus
vanishing forces every coefficient to vanish.  Substitution on the
generators proves the multiplication identity.  On monomials,
\[
       B_F(X_eP,Q)=B_F(P,\partial_eQ).
 \tag{JIO.4}
\]
The chain rule gives
\(\partial_eL(A^*)Q=\sum_f\overline{A_{fe}}L(A^*)\partial_fQ\).
Induct on the degree of \(P\).  The constant case uses preservation
of the constant coefficient.  For \(P=X_eP_0\), expand \(L(A)X_e\),
move each \(X_f\) across the form using (JIO.4), and apply the
induction hypothesis to \(P_0\) and \(\partial_fQ\).  The resulting
sum is exactly the chain rule just displayed.  Addition and scalar
multiplication finish the induction.
% END JIO1

% BEGIN JIO2
\paragraph{JIO2: both original Fock actions are holomorphic product representations.}
\label{claim:JIO2}
For \(k_\tau\in K_{\epsilon_\tau}^{\mathbb C}\) and
\(u_\tau\in\operatorname{GL}_n(\mathbb C)\), define the actual
substitution by
\[
 X_{\tau,i,j}\longmapsto
 \begin{cases}
 \displaystyle\sum_{a,b}(k_\tau)_{ai}(u_\tau^{-1})_{jb}X_{\tau,a,b},
                                           &\epsilon_\tau(i)=+,\\
 \displaystyle\sum_{a,b}(k_\tau^{-1})_{ia}(u_\tau)_{bj}X_{\tau,a,b},
                                           &\epsilon_\tau(i)=-.
 \end{cases}
 \tag{JIO.5}
\]
Call it \(f_{\rm sub}(k,u)\).  Put \(P_\pm=(1\pm J_\epsilon)/2\),
\(d_\pm(k)=\det(P_\pm k+P_\mp)\), and, for arbitrary integer-valued
\(x_1,x_2,x_3\), put
\[
 \begin{split}
 c(k,u)&=\prod_\tau d_+(k_\tau)^{x_{1,\tau}}
                      d_-(k_\tau)^{x_{2,\tau}}
                      \det(u_\tau)^{x_{3,\tau}},\\
 F(k,u)&=c(k,u)f_{\rm sub}(k,u),\\
 F_K(g)&=F(g^{-\mathsf t},1),\qquad F_U(u)=F(1,u).
 \end{split}
 \tag{JIO.6}
\]
These are the source's \texttt{fchar}, \texttt{fact},
\texttt{factK}, \texttt{factU}.  Both \(F_K\) and \(F_U\) satisfy
\texttt{IsHolFam} for \(B_F\), with sign families respectively
\(\epsilon\) and the all-positive sign.  Moreover \(F_U(u)\)
commutes with \(F_K(g)\) for every allowed \(u,g\).

\emph{Proof.}
Commutation with \(J_\epsilon\) says \(k_{ai}=0\) whenever the
two signs differ, and gives the same assertion for \(k^{-1}\).
Consequently the two cases in (JIO.5) compose separately.
In the positive block the coefficients of the composite are
\((kk')_{ai}((uu')^{-1})_{jb}\); in the negative block they are
\(((kk')^{-1})_{ia}(uu')_{bj}\).  Thus \(f_{\rm sub}\) is the
linear substitution \(L(A(k,u))\) of JIO1 and
\[
 A(kk',uu')=A(k,u)A(k',u'),\qquad
 A(k,u)^*=A(k^*,u^*).
 \tag{JIO.7}
\]
For the second equality, for example, the conjugate transpose of
the positive coefficient \(k_{ai}(u^{-1})_{jb}\) is
\((k^*)_{ai}((u^*)^{-1})_{jb}\); the negative coefficient is checked
in the same order.

In a basis ordered by the signs, \(d_+(k)\) and \(d_-(k)\) are the
determinants of the two diagonal blocks.  They are nonzero, multiply
under multiplication, conjugate under \(k\mapsto k^*\), and satisfy
\(d_+(k)d_-(k)=\det k\).  Hence JIO1 proves
\[
 F(kk',uu')=F(k,u)F(k',u'),\qquad
 B_F(F(k,u)P,Q)=B_F(P,F(k^*,u^*)Q).
 \tag{JIO.8}
\]
The identity action is immediate.  Inverse transpose is a group
automorphism of the complex block group and commutes with \(^{*}\).
It follows that both restrictions in (JIO.6) have the required
representation and adjoint identities.  The two factor groups
commute, proving the last assertion.

For any two fixed polynomials, all relevant coefficients are finite
sums of products of matrix entries, inverse entries and integer
powers of the three nonzero determinants.  The formula
\(A^{-1}=\operatorname{adj}(A)/\det A\) proves that they are
holomorphic on the prescribed block groups.  If the source's
ambient-matrix formulation of differentiability is used, these
same formulas are holomorphic on a neighborhood of each block
matrix where the two block determinants and \(\det u\) remain
nonzero.  Their restrictions give the stated differentiability.
Finally JIO1 supplies the anisotropy condition.  This proves the
complete two-family contract.
% END JIO2

% BEGIN JIO3
\paragraph{JIO3: the original minor product and its triangular eigenvalues.}
\label{claim:JIO3}
Fix \(p+q=m\le n\), with the first \(p\) signs positive, and
nonnegative sequences \(a^+_0\ge\cdots\ge a^+_{p-1}\),
\(a^-_0\ge\cdots\ge a^-_{q-1}\).  Set \(a^+_p=a^-_q=0\).
At one place define
\[
 \begin{split}
 D^+_s&=\det(X_{\beta,\alpha})_{0\le\alpha,\beta<s},\\
 D^-_s&=\det(X_{p+\beta,n-s+\alpha})_{0\le\alpha,\beta<s},\\
 \Delta_\tau&=\prod_{s=1}^{p}(D^+_s)^{a^+_{s-1}-a^+_s}
               \prod_{s=1}^{q}(D^-_s)^{a^-_{s-1}-a^-_s}.
 \end{split}
 \tag{JIO.9}
\]
The product \(\Delta=\prod_\tau\Delta_\tau\) is exactly
\texttt{delta p q ap am}, and is nonzero.
If \(k\) is upper triangular on the positive block, lower triangular
on the negative block, and \(u\) is lower triangular, then
\[
 \begin{split}
 f_{\rm sub}(k,u)D^+_s
    &=\left(\prod_{a<s} k_{aa}u_{aa}^{-1}\right)D^+_s,\\
 f_{\rm sub}(k,u)D^-_s
    &=\left(\prod_{a<s} k_{p+a,p+a}^{-1}
                                  u_{n-1-a,n-1-a}\right)D^-_s.
 \end{split}
 \tag{JIO.10}
\]
The resulting target-row exponents of \(\Delta_\tau\) are
\(a^+_a\) on the positive block and \(-a^-_{a-p}\) on the
negative block; its source-column exponents are
\(-a^+_j\) for \(j<p\), \(+a^-_{n-1-j}\) for \(n-1-j<q\),
and zero elsewhere.

\emph{Proof.}
For \(D^+_s\), upper triangularity of \(k\) makes the first \(s\)
target rows stable under the coefficients \(k_{ai}\);
lower triangularity of \(u^{-1}\) makes the first \(s\)
source columns stable.  The transformed square matrix is the
product of the corresponding triangular submatrices and the
original minor matrix.  Its determinant gives the first formula.
For \(D^-_s\), the first \(s\) rows of the negative block are
stable under \(k^{-1}\), and the last \(s\) columns are stable
under the coefficients \(u_{bj}\).  The same determinant
multiplication proves the second formula, with the displayed
last-column order.  Transposing a displayed minor does not
change its determinant.

For any nonzero \(z_0,\ldots,z_{r-1}\) and decreasing
nonnegative \(a_0,\ldots,a_{r-1}\), the exponent of \(z_j\) in
\[
          \prod_{s=1}^{r}
                (z_0\cdots z_{s-1})^{a_{s-1}-a_s}
                  =\prod_{j=0}^{r-1}z_j^{a_j}
 \tag{JIO.11}
\]
is the telescoping sum
\(\sum_{s=j+1}^{r}(a_{s-1}-a_s)=a_j\).
Apply this to (JIO.10), including inverses.
All exponents in (JIO.9) are natural numbers by the assumed
antitonicity, so the source's truncated natural subtraction
agrees with ordinary subtraction.

Every determinant in (JIO.9) is a nonzero polynomial: its diagonal
monomial has coefficient \(1\) and occurs in no other permutation
term.  Their product is nonzero in the polynomial domain.  This
also covers \(p=0\), \(q=0\) and the empty products.  The inequality
\(p+q\le n\) ensures that all specified minors have their original
indices and sizes; no overlap or index truncation has been assumed
away.
% END JIO3

% BEGIN JIO4
\paragraph{JIO4: the negative-tail reversal gives the exact joint bottom vector.}
\label{claim:JIO4}
At each place let \(\kappa_\tau\) be the permutation matrix fixing
the first \(p_\tau\) indices and reversing the last \(q_\tau\).
It is real, unitary, sign-preserving and satisfies \(\kappa_\tau^2=1\).
Put \(\Delta'=F(\kappa,1)\Delta\).  Then
\[
 F(\kappa,1)\Delta'=\Delta,\quad \Delta'\ne0,
 \quad
 \begin{cases}
 \Delta'\text{ is bottom for }F_K\text{ of weight }-\mathrm{wtK},\\
 \Delta'\text{ is bottom for }F_U\text{ of weight }\mathrm{wtU},
 \end{cases}
 \tag{JIO.12}
\]
where the literal integral weights are
\[
 \begin{split}
 \mathrm{wtK}_a&=
   \begin{cases}a^+_a+x_1,&a<p,\\
                -a^-_{m-1-a}+x_2,&a\ge p,\end{cases}\\
 \mathrm{wtU}_j&=-[j<p]\,a^+_j
                  +[n-1-j<q]\,a^-_{n-1-j}+x_3 .
 \end{split}
 \tag{JIO.13}
\]

\emph{Proof.}
If \(k\) is upper triangular in both sign blocks,
\(\kappa k\kappa\) is upper triangular on the positive block
and lower triangular on the negative block.  JIO2 and JIO3 give
\[
 F(k,u)\Delta'
       =\left(\prod_a k_{aa}^{\mathrm{wtK}_a}
                    \prod_j u_{jj}^{\mathrm{wtU}_j}\right)\Delta'
 \tag{JIO.14}
\]
when \(u\) is lower triangular, with products also over places.
Indeed the reversal changes the negative row exponent from
\(-a^-_{a-p}\) to \(-a^-_{m-1-a}\).  The determinant characters
are unchanged by this conjugation, and contribute exactly
\(x_1,x_2,x_3\).

For a lower-unitriangular \(g\), the matrix \(g^{-\mathsf t}\)
is upper-unitriangular in the same blocks.  Formula (JIO.14)
therefore gives \(F_K(g)\Delta'=\Delta'\).  On
\(g=\operatorname{diag}(e^{i\theta_a})\), inverse transpose changes
the angles to \(-\theta_a\), giving the weight \(-\mathrm{wtK}\).
The \(F_U\) statement follows directly from lower-unitriangular
and diagonal \(u\) in (JIO.14).  One-place specialization gives
the full \texttt{IsBotFam} predicates.  Finally JIO2 implies
\(F(\kappa,1)^2=1\); hence both the equality and nonvanishing in
(JIO.12) follow from JIO3.  The determinant of the reversal may
be \(-1\), but its integer powers occur twice in the square, so
no unrecorded sign remains.
% END JIO4

% BEGIN JIO5
\paragraph{JIO5: an explicit holomorphic extension of the compact action.}
\label{claim:JIO5}
Let \(\pi:K_\epsilon\to\operatorname{End}_{\mathbb C}(S)\)
be a continuous unitary representation on a finite-dimensional
Hilbert space.  There is a holomorphic matrix-valued function
\(\tau\) on \(\operatorname{GL}_m(\mathbb C)\), equal to \(\pi\)
on \(K_\epsilon\), whose restriction to \(K_\epsilon^{\mathbb C}\)
is a representation and satisfies \(\tau(g)^*=\tau(g^*)\).
Thus this restriction is an \texttt{IsHolRep}.

\emph{Proof.}
Use probability Haar measure on the compact block group and put
\[
 p(k)=1+2m+2\Re\operatorname{tr}k,\quad
 c_N=\int_Kp(k)^N\,dk,\quad
 A_N=\int_Kp(k)^N\pi(k)\,dk .
 \tag{JIO.15}
\]
For unitary \(k\),
\(\operatorname{tr}((k+1)^*(k+1))=2m+2\Re\operatorname{tr}k\)
is nonnegative.  Hence \(p(k)\ge1\).  Also
\(p(1)-p(k)=\operatorname{tr}((k-1)^*(k-1))\), so \(1\)
is its unique maximum.  The probability measures
\(c_N^{-1}p(k)^N\,dk\) concentrate at \(1\).
For completeness, outside any neighborhood \(U\) of \(1\) the
maximum is \(r<p(1)\); choose \(r<r'<p(1)\) and an open
neighborhood \(V\) on which \(p>r'\).  Haar measure of \(V\)
is positive (finitely many translates cover \(K\)), and the
mass outside \(U\) is at most
\(\mu(V)^{-1}(r/r')^N\).  This tends to zero.
Uniform continuity of \(\pi\) gives \(c_N^{-1}A_N\to1\) in
operator norm.  Choose \(N\) so close that its distance to \(1\)
is less than \(1\).  The convergent Neumann series makes \(A_N\)
invertible.

Define, for \(g\in\operatorname{GL}_m(\mathbb C)\),
\[
 \tau(g)=
 \left(\int_K
 [1+2m+\operatorname{tr}(g^{-1}k)+\operatorname{tr}(k^*g)]^N
                       \pi(k)\,dk\right)A_N^{-1}.
 \tag{JIO.16}
\]
Expanding the finite power shows directly that this is a finite
sum of holomorphic functions of entries of \(g,g^{-1}\), with
constant operator coefficients.  For \(g\in K\) the bracket is
\(p(g^{-1}k)\); left invariance of Haar measure and the group law
give the integral \(\pi(g)A_N\).  Thus \(\tau(g)=\pi(g)\).

We need the following uniqueness fact, including its justification.
A holomorphic function on the complex block group which vanishes
on \(K\) vanishes everywhere.  For \(y\) in the complex block
group, take the positive Hermitian logarithm
\(H=\log(y^*y)\), obtained by diagonalizing its two positive
Hermitian blocks.  It commutes with \(J_\epsilon\), and
\(u=y e^{-H/2}\) belongs to \(K\).  The entire curve
\(z\mapsto u e^{zH}\) lies in the complex block group and lies
in \(K\) when \(z\) is imaginary.  The one-variable identity
theorem, applied to matrix entries, gives vanishing at
\(z=1/2\), namely at \(y\).

Apply this uniqueness fact first in one variable and then the
other to
\(\tau(ab)-\tau(a)\tau(b)\).  Its vanishing on \(K\times K\)
extends to the whole product of complex block groups.
The same fact applied entrywise to
\(\tau(g)-\tau(g^*)^*\), a holomorphic function because the two
conjugations cancel, extends the adjoint identity from \(K\).
The original positive inner product is anisotropic, completing
the claimed contract.  The construction uses no unproved
complexification or algebraic-representation existence assertion.
% END JIO5

% BEGIN JIO6
\paragraph{JIO6: the conjugated compact type has its required bottom vector.}
\label{claim:JIO6}
Let \(S\) satisfy the original
\(\mathrm{IsKTypeOf}(\epsilon,\sigma,\lambda,S)\): it is finite
dimensional, compact-stable and irreducible, has a nonzero
weight-\(\lambda\) vector, and every occurring weight \(\mu\)
satisfies the source's compact dominance
\(\lambda-\mu=\sum_{a<b,\epsilon_a=\epsilon_b}c_{ab}(e_a-e_b)\)
with \(c_{ab}\in\mathbb N\).
Choose a unit vector \(w\in S\) of weight \(\lambda\).
There is an \texttt{IsHolRep} \(\tau\) on \(S\) such that
\[
       \tau(k)s=\sigma(\overline{k})s\quad(k\in K_\epsilon),
       \qquad
       w\text{ is bottom of weight }-\lambda.
 \tag{JIO.17}
\]

\emph{Proof.}
Complex conjugation preserves the compact block group and its
group law.  Apply JIO5 to \(k\mapsto\sigma(\overline{k})|_S\).
Its compact invariant subspaces are exactly those of \(\sigma\),
so irreducibility is unchanged.  On the diagonal torus,
conjugation reverses all angles, giving weight \(-\lambda\)
for \(w\).

Let \(D\) be the complex derivative of \(\tau\) at \(1\).
For \(a>b\) in the same sign block, differentiating
\[
 t(1+zE_{ab})t^{-1}
      =1+e^{i(\theta_a-\theta_b)}zE_{ab},
 \qquad t=\operatorname{diag}(e^{i\theta_j}),
 \tag{JIO.18}
\]
shows that \(D(E_{ab})w\), if nonzero, has \(\tau\)-weight
\(-\lambda+e_a-e_b\).  Its \(\sigma\)-weight is therefore
\(\lambda-(e_a-e_b)\).  Compact dominance would express
\(e_a-e_b\) as a nonnegative integral combination of roots
\(e_c-e_d\) with \(c<d\) in equal-sign blocks.
Apply the functional \(v\mapsto\sum_j jv_j\):
the left side is \(a-b>0\), whereas every summand on the right
is nonpositive.  This is impossible.  Thus \(D(E_{ab})w=0\).

For the one-parameter subgroup \(1+zE_{ab}\), \(E_{ab}^2=0\).
Differentiating its group law at any \(z\) now gives derivative
\(\tau(1+zE_{ab})D(E_{ab})w=0\).  Hence it fixes \(w\)
for every complex \(z\).  Ordinary elimination, starting with
the bottom row and moving upward, writes every lower-unitriangular
matrix in each sign block as a product of these elementary
matrices.  They all fix \(w\), so their product fixes \(w\).
This proves the full bottom condition in (JIO.17), including
all lower-unitriangular matrices rather than only their tangent
directions.
% END JIO6

% BEGIN JIO7
\paragraph{JIO7: a bottom matrix coefficient depends only on its weight.}
\label{claim:JIO7}
Let \((P,B,\pi)\) and \((P',B',\pi')\) be two
\texttt{IsHolRep}s for the same sign group.  If \(x,y\) and
\(x',y'\) are bottom vectors of the same real weight \(\mu\), then
for every \(k\in K_\epsilon\),
\[
 B(x,\pi(k)y)\,B'(x',y')
             =B(x,y)\,B'(x',\pi'(k)y').
 \tag{JIO.19}
\]
No finite-dimensionality of \(P\) or \(P'\) is required.

\emph{Proof.}
First the unitary diagonal law extends to positive diagonal
matrices:
\[
 \pi(\operatorname{diag}(e^{s_a}))y
                         =e^{\sum_a\mu_as_a}y,\quad s_a\in\mathbb R.
 \tag{JIO.20}
\]
Indeed take scalar coefficients against any fixed \(z\) along
\(\operatorname{diag}(e^{i\zeta s_a})\).
Both sides are entire in \(\zeta\) and agree for real \(\zeta\).
They agree at \(\zeta=-i\).  If their vector difference is \(v\),
test with \(z=v\); anisotropy of \(B\) makes the difference zero.
The same proof works for every bottom vector and for \(B'\).

A lower triangular block matrix \(b\) with positive diagonal
is the product of its positive diagonal part and a lower
unipotent.  Thus it acts on every bottom vector by the same
positive number
\(r(b,\mu)=\exp(\sum_a\mu_a\log b_{aa})\).
For any \(g\in K_\epsilon^{\mathbb C}\), Gram--Schmidt on its
columns in reverse order, separately in the two blocks, gives
\(g=kb\) with \(k\) unitary and \(b\) lower triangular with
positive diagonal.  At each step the vector to be normalized
is nonzero because the remaining columns are independent.
Adjointness and the group law then give
\[
             B(x,\pi(g^*g)y)=r(b,\mu)^2B(x,y).
 \tag{JIO.21}
\]
The same scalar appears for \(B',\pi'\).

Finally diagonalize the two unitary blocks of \(k\).  Choosing
real arguments of their eigenvalues gives \(k=e^{iH}\) with
\(H=H^*\) commuting with \(J_\epsilon\).
Along \(e^{zH}\), the difference between the two sides of
(JIO.19) is entire.  For real \(z=t\),
\(e^{tH}=(e^{tH/2})^*e^{tH/2}\), so (JIO.21) makes that
difference zero.  The identity theorem gives zero at \(z=i\),
which is exactly (JIO.19).  The proof uses neither division by
\(B(x,y)\) nor any assertion that this coefficient is nonzero.
% END JIO7

% BEGIN JIO8
\paragraph{JIO8: the product orbit map is constructed from its Gram relations.}
\label{claim:JIO8}
Let two \texttt{IsHolFam}s have the same sign family and
bottom weights \(\mu_\tau\).  If \(x\in P\), \(x'\in P'\) are
common bottom vectors and \(x'\ne0\), there exists a linear
\(\Psi:P'\to P\) such that for every compact tuple \(u\),
\[
                    \Psi(\pi'(u)x')=\pi(u)x.
 \tag{JIO.22}
\]
The assertion is the exact \texttt{exists\_orbit\_map} contract;
it does not claim equivariance on an arbitrary complement.

\emph{Proof.}
Since \(B'(x',x')\ne0\), JIO7 gives, for each one-place compact
element \(k\), a scalar
\[
        c_\tau(k)=
          \frac{B'(x',\pi'(\operatorname{single}_\tau k)x')}
               {B'(x',x')}
 \tag{JIO.23}
\]
which computes the coefficient on any pair of bottom vectors
at that place, for either representation.
Operators at distinct places commute.  Applying an operator at
one place preserves all bottom laws at the other places.
The one-place adjoint identities also give the adjoint identity
for a whole tuple by writing it as the finite product of its
one-place tuples.
Successively apply JIO7 to obtain the product coefficient
\(\prod_\tau c_\tau(u_\tau)\).
Adjointness gives the resulting orbit Gram identity
\[
 \begin{split}
 &B\!\left(\sum_j a_j\pi(u_j)x,\sum_l b_l\pi(v_l)x\right)
                         B'(x',x')\\
 &\hspace{5mm}=
 B(x,x)\,
 B'\!\left(\sum_j a_j\pi'(u_j)x',
                         \sum_l b_l\pi'(v_l)x'\right).
 \end{split}
 \tag{JIO.24}
\]
Here the one-place argument is applied to \(u_j^{-1}v_l\);
the inverse is its conjugate transpose.

If a finite linear relation between source orbit vectors is zero,
use the same relation on both sides of (JIO.24).  The target
vector has zero \(B\)-square and therefore is zero.  Thus (JIO.22)
defines a well-defined linear map on the span of the source
orbit.  Extend a basis of that span to a basis of \(P'\) and put
the map equal to zero on the added basis vectors.  This is the
required linear map.  If \(x=0\), the same argument simply gives
the zero map on the orbit span.
% END JIO8

% BEGIN JIO9
\paragraph{JIO9: a universal finite-type operator, with its operator-span control.}
\label{claim:JIO9}
Let \(S,\sigma,\lambda\) be as in JIO6 and let
\((P,B,\pi)\) be any \texttt{IsHolRep} for the same signs.
There are a unit vector \(w\in S\) and a linear map
\(T:S\to\operatorname{End}_{\mathbb C}(P)\) such that
\[
 \begin{gathered}
 T(s)\in\operatorname{span}_{\mathbb C}\{\pi(k):k\in K_\epsilon\},
 \\
 T(w)y=y,\qquad
 T(\sigma(k)s)y=\pi(\overline{k})T(s)y
 \quad\text{for every bottom }y\text{ of weight }-\lambda .
 \end{gathered}
 \tag{JIO.25}
\]
The same \(T\) works for every such \(y\).

\emph{Proof.}
Use the holomorphic representation \(\tau\) of JIO6.
On the free vector space of finitely supported complex functions
on \(K_\epsilon\), define
\[
       \Phi(a)=\sum_k a_k\tau(k)w,\qquad
       \Upsilon(a)=\sum_k a_k\pi(k).
 \tag{JIO.26}
\]
The range of \(\Phi\) is nonzero and compact-stable.  Irreducibility
therefore makes it all of \(S\).  Choose a linear right inverse
\(R\) by choosing preimages of a basis of \(S\), and set
\(T=\Upsilon R\).  This immediately gives the operator-span
assertion.

For every bottom \(y\), JIO7 and \(\|w\|=1\) give
\[
 B(\Upsilon(a)y,\Upsilon(b)y)
                      =B(y,y)\langle\Phi(a),\Phi(b)\rangle .
 \tag{JIO.27}
\]
Thus \(\Phi(a)=0\) implies \(\Upsilon(a)y=0\), by anisotropy.
Consequently
\(\Upsilon(R\Phi(a))y=\Upsilon(a)y\), simultaneously for all
bottom \(y\).
Take \(a\) to be the identity basis vector to get \(T(w)y=y\).
Left translation of \(a\) by \(k\) gives
\(T(\tau(k)s)y=\pi(k)T(s)y\).
Replace \(k\) by \(\overline{k}\) and use
\(\tau(\overline{k})s=\sigma(k)s\), obtaining (JIO.25).
The chosen right inverse depends on \(S,\tau,w\), not on \(y\),
which is the needed simultaneous quantifier.
% END JIO9

% BEGIN JIO10
\paragraph{JIO10: the operator acts on the whole original \(\mathrm{LowT}(M)\).}
\label{claim:JIO10}
For any original \(\mathrm{DSModel}\ M\) and any
\texttt{IsHolFam} \((P,B,\varpi)\) with signs \(M.\epsilon\),
there is
\(\mathfrak m:\mathrm{LowT}(M)\to\operatorname{End}_{\mathbb C}(P)\)
with all three conclusions of \texttt{lowest\_types\_act}:
it commutes with every operator commuting with all compact
\(\varpi\)-operators; for every compact tuple \(\kappa\) there are
nonzero \(v_\tau\in\mathrm{lowT}(M,\tau)\) with
\[
 \mathfrak m(\otimes_\tau v_\tau)y
                  =\varpi(\overline{\kappa})y
       \quad\text{for every common bottom }y
                \text{ of weight }-\mathrm{lowK}(M);
 \tag{JIO.28}
\]
and replacement of \(v_\tau\) by \(\sigma_\tau(k)v_\tau\)
changes the value to
\(\varpi(\operatorname{single}_\tau\overline{k})
          \mathfrak m(\otimes v)y\).

\emph{Proof.}
HCU8 gives, at every place, a finite-dimensional smooth subspace
\(S_\tau\) satisfying the original lowest \(\mathrm{IsKTypeOf}\).
It is contained in the original \(\mathrm{lowKSub}\), since the
latter is the supremum of \emph{all} subspaces of that type.
Let \(P_\tau\) be the orthogonal projection from \(M.H_\tau\)
to \(S_\tau\).  It commutes with \(\sigma_\tau(k)\):
unitarity preserves both \(S_\tau\) and its orthogonal complement,
so the two components of a vector are carried to the corresponding
two components of its translate.

Apply JIO9 to \(S_\tau\) and the one-place representation.
Precompose its \(T_\tau\) with \(P_\tau\) and with the literal
inclusion
\(\mathrm{lowT}(M,\tau)\subset M.\mathrm{sm}_\tau\subset M.H_\tau\).
Call the resulting operator-valued map \(U_\tau\).
Each \(U_\tau(v)\) belongs to the span of one-place compact
\(\varpi\)-operators.  Hence \(U_\tau(v)\) and \(U_\rho(w)\)
commute if \(\tau\ne\rho\), by the product-group law and
bilinearity of commutation.  Therefore
\[
          \mathfrak m(\otimes_\tau v_\tau)
                           =\prod_\tau U_\tau(v_\tau)
 \tag{JIO.29}
\]
is multilinear and defines the asserted linear map on the
algebraic tensor product.  The product is independent of the
chosen ordering because its factors commute.
The operator-span assertion also proves the commutant conclusion.

Choose \(v_\tau=\sigma_\tau(\kappa_\tau)w_\tau\) from JIO9,
regarded as an element of the original \(\mathrm{lowT}\).
It lies there because \(S_\tau\) is smooth, compact-stable and
contained in \(\mathrm{lowKSub}\); its norm is \(1\), so it is
nonzero.  JIO9 gives
\(U_\tau(v_\tau)y=
 \varpi(\operatorname{single}_\tau\overline{\kappa_\tau})y\).
Multiplying these equalities, commuting the other-place factors
at each step, gives (JIO.28).
For an arbitrary updated tuple, projection equivariance and
the last equality in (JIO.25) give the one-place replacement
law; the remaining factors commute past that one-place operator.
This proves the full original contract on \(\mathrm{LowT}(M)\).
It has not replaced that space by \(S_\tau\), nor required
HolShape or multiplicity one.
% END JIO10

% BEGIN JIO11
\paragraph{JIO11: the literal coefficient polynomial is a nonzero bottom vector.}
\label{claim:JIO11}
For original dominant integral weights \(\lambda_\tau\), put
\(d_\tau=\lambda_\tau.\mathrm{detShift}\) and
\(\nu_{\tau,j}=\lambda_{\tau,j}-d_\tau\in\mathbb N\).
On the source's square-matrix polynomial space define
\[
 V_\lambda(u)=
       \left(\prod_\tau\det(u_\tau)^{-d_\tau}\right)
                    \operatorname{blockAct}(u^{-1}),\qquad
 e_\lambda=\prod_\tau
            \prod_{s=1}^{n}
        \det(X_{\tau,a,b})_{a,b<s}^{\,\nu_{\tau,s-1}-\nu_{\tau,s}},
 \quad \nu_{\tau,n}=0.
 \tag{JIO.30}
\]
Then \(V_\lambda\) is an \texttt{IsHolFam} for the all-positive
sign and the factorial form.  The vector \(e_\lambda\) is nonzero,
belongs to the literal \(\mathrm{VV}(\mathbb C,\lambda)\), and
is a common bottom vector of weight \(-\lambda\).
These are precisely the four conclusions of
\texttt{vact\_isHolFam\_isBotFam}.

\emph{Proof.}
The last entry of a decreasing integral weight is its minimum,
so \(\nu_{\tau,j}\ge0\) and is decreasing; for \(n=0\) the
source chooses \(d_\tau=0\) and all products are empty.
The displayed polynomial is exactly the source's tensor of
the \texttt{extremal} polynomials.

The substitution \(X_{ij}\mapsto\sum_a(u^{-1})_{ia}X_{aj}\)
has the group law in this order.  JIO1, now with all rows as
variables, gives its adjoint under \(u\mapsto u^*\).
The determinant twist has the same adjoint rule.
Inverse entries and integer determinant powers are holomorphic
on the invertible matrices, proving all family conditions as
in JIO2.

For a lower triangular \(u\), its inverse is lower triangular;
the first \(s\) rows of the leading minor are stable.  The
minor scales by \(\prod_{a<s}u_{aa}^{-1}\).
JIO3's telescoping identity gives row exponent \(-\nu_j\);
the determinant twist adds \(-d\), giving \(-\lambda_j\).
Thus lower unipotents fix \(e_\lambda\), and the diagonal
unitary character is the required one.

For membership, the source defines \(\mathrm{dualWeyl}(\nu)\)
as the column-multihomogeneous polynomials of degree \(\nu\)
fixed by every operation adding an earlier column to a later
one, even with a formal indeterminate as coefficient.
Every leading determinant is fixed by such an operation:
either the later column is absent, or both columns are present
and determinant alternation cancels the new term.
Telescoping gives column degree \(\nu\).
Consequently each factor lies in its literal
\(\mathrm{nablaSq}\), and their tensor product lies in the
defining span \(\mathrm{VV}\).
Evaluating every square matrix at the identity makes every
leading determinant \(1\), so \(e_\lambda\ne0\).
This establishes the complete original polynomial contract.
% END JIO11

% BEGIN JIO12
\paragraph{JIO12: the parameter-free map and its exact unit-vector normalization.}
\label{claim:JIO12}
Take all the original inputs of \texttt{exists\_io\_poly}:
\(D,V,W,\theta,M,c,\lambda,c_0,B_S,p,q,a^+,a^-,x_1,x_2,x_3\);
\(\theta\) is an \(\mathrm{IsInfFam}\), \(c\) satisfies \(c.\mathrm{Is}\),
and the same \(c_0\) satisfies
\(\mathrm{Bridge.CoefMapTie}(V,c,c_0,\rho_{V,r})\).
Assume \(B_{S,\tau}\) is invertible, \(m\le n\), \(p_\tau+q_\tau=m\),
the ordered sign condition, both antitonicity conditions of JIO3,
and the exact two weight identities
\[
       \mathrm{lowK}(M)_{\tau,a}=\mathrm{wtK}_{\tau,a},
              \qquad \mathrm{wtU}_{\tau,j}=-\lambda_{\tau,j}.
 \tag{JIO.31}
\]
Then there is one bilinear map
\[
           \mathrm{io}:\mathrm{LowT}(M)\times c.W
                                  \longrightarrow\mathcal P
 \tag{JIO.32}
\]
satisfying both polynomial covariance laws in the source
\texttt{exists\_io\_poly}, and taking one tuple of unit lowest
vectors and one unit coefficient vector to the \emph{exact}
polynomial \(\Delta\).

\emph{Proof.}
Use CFB5 with this same \(c_0,B_S\).
It supplies a linear map \(j:c.W\to\operatorname{Pol}\) with
\(\operatorname{range}j=\mathrm{VV}(\mathbb C,\lambda)\) and
\[
 j(c.\rho(u)w)=V_\lambda(u_B)j(w),\qquad
 (u_B)_\tau=B_{S,\tau}^{-1}
                     \operatorname{compAt}_\tau(\mathrm{mt}(u))
                       B_{S,\tau}
 \tag{JIO.33}
\]
for every original \(u\in G_\infty(V)\).
CFB6 also proves that the compact \(V_\lambda\)-orbit of
\(e_\lambda\) spans this exact \(\mathrm{VV}\).
These are explicit proved coefficient interfaces, not an
unspecified intertwiner premise.

JIO2, JIO4, JIO11 and the second identity in (JIO.31) allow
JIO8 to construct
\(\Psi:\operatorname{Pol}\to\mathcal P\) with
\[
             \Psi(V_\lambda(u)e_\lambda)=F_U(u)\Delta'
                      \quad(u\text{ compact}).
 \tag{JIO.34}
\]
On \(\mathrm{VV}\) this implies two facts:
\(\Psi(F)\) is bottom for \(F_K\) of weight
\(-\mathrm{lowK}(M)\), and
\(\Psi(V_\lambda(u)F)=F_U(u)\Psi(F)\) for compact \(u\).
For orbit generators the first follows from commutation of
\(F_U\) and \(F_K\), and the second from the group law.
Both properties extend to their span by linearity.
No assertion is made about equivariance on the chosen complement.

Let \(\mathfrak m\) be JIO10 for \(F_K\).
Choose its nonzero tuple \(v^0\) for the reversal \(\kappa\).
Choose \(w^0\) with \(j(w^0)=e_\lambda\); it is nonzero because
\(e_\lambda\ne0\).  Equation (JIO.34) at \(u=1\) gives
\(\Psi(jw^0)=\Delta'\).
Write the positive numbers
\[
 N_\tau=\|v^0_\tau\|,\quad N_0=\|w^0\|,\quad
 C=N_0\prod_\tau N_\tau,\qquad
          \mathrm{io}(x,w)=C\,\mathfrak m(x)\Psi(jw).
 \tag{JIO.35}
\]
All norms here are the original Hilbert and coefficient norms.
They are nonzero, and the expression is complex-bilinear.

If \(u\in G_\infty(V)\) and all \(u_B\) are unitary,
(JIO.33)--(JIO.34) and the commutant property of \(\mathfrak m\)
give
\[
                  F(1,u_B)\mathrm{io}(x,w)
                             =\mathrm{io}(x,c.\rho(u)w).
 \tag{JIO.36}
\]
For \(k\in K_{\epsilon_\tau}\) and an original updated lowest
vector \(v'_\tau=\sigma_\tau(k)v_\tau\), JIO10 gives the
operator \(F_K(\operatorname{single}_\tau\overline{k})\).
Unitarity implies
\(((\overline{k})^{-1})^{\mathsf t}=k\); therefore
\[
 F(\operatorname{single}_\tau k,1)\mathrm{io}(\otimes v,w)
                 =\mathrm{io}(\otimes(v[\tau\mapsto v'_\tau]),w).
 \tag{JIO.37}
\]
This checks the otherwise delicate inverse-transpose convention.

Finally put \(v_\tau=N_\tau^{-1}v^0_\tau\) and
\(w=N_0^{-1}w^0\).  Their norms are \(1\).
Tensor multilinearity and (JIO.28) show
\[
 \begin{split}
 \mathrm{io}(\otimes v,w)
   &=C\left(\prod_\tau N_\tau^{-1}\right)N_0^{-1}
                         F_K(\overline{\kappa})\Delta'\\
   &=F(\kappa,1)\Delta'=\Delta.
 \end{split}
 \tag{JIO.38}
\]
The scalar cancels to \(1\), with no phase or unspecified
normalization factor.  This proves the entire parameter-free
source conclusion.
% END JIO12

% BEGIN JIO13
\paragraph{JIO13: the stored place-model value is an actual theta vacuum.}
\label{claim:JIO13}
For the original standard character and infinity family, suppose
\(\mathrm{IsThetaAt}(\psi,e,\Lambda,c,f_0,\Phi,f)\) holds and
\(f_0\) is an original theta function.  Then \(f\) is an original
theta function.  Consequently the original
\(R.\mathrm{thP}_{vP}\), together with \(R.\mathrm{IsModels}\),
gives, for every \(i,t\), the theta property of \(R.vP(i,t)\).
The stored vacuum reading and the original creators then give
\[
 \begin{gathered}
 f_{i,t}=R.vP(i,t)\text{ is Fock-vacuum at every real place},\\
 \mathrm{pt}_{i,t}(P)=
 \mathrm{pureTensorO}\bigl(\psi,\theta,V,W,1,\mathrm{cart}(M),
                    R.BV(i),\mathrm{onb}(M),f_{i,t},P\bigr)\\
 \text{ is a theta function for every }P .
 \end{gathered}
 \tag{JIO.39}
\]
Here \(M=R.MW(i,\mathrm{sg0})\) and all objects are the stored ones.

\emph{Proof.}
Let \(\Theta\) be the vector subspace of original theta functions.
It is a subspace because sums use the intersection of their open
finite levels, and the continuity, quasi-periodicity and smoothness
conditions are linear.  Every original Heisenberg operator \(H_w\)
preserves it: use the Heisenberg multiplication formula for
quasi-periodicity, and shrink the finite level to the kernel of
the extra pairing character with the fixed finite part of \(w\).
This is an open level: on a sufficiently small compact open
additive subgroup, continuity puts the character in an arc of
the circle containing no nontrivial subgroup, so its image is
trivial.  Translating the defining real orbit proves smoothness.
This is also the theta-space calculation of CW12.

Let \(\ell\) be any algebraic linear functional annihilating
\(\Theta\), and define a linear operator on all functions by
\[
                       (Ag)(w)=\ell(H_wg).
 \tag{JIO.40}
\]
The Heisenberg cocycle identity says
\(AH_u=H_uA\): at \(w\), both sides are the same phase times
\(\ell(H_{w+u}g)\).  Since \(H_wf_0\in\Theta\), \(Af_0=0\).
Choose a place model \(T_0\) supplied by the existential part of
\(\mathrm{IsThetaAt}\).  The map \(T_0+AT_0\) is another place
model for the identical \(e,\Lambda,c,f_0\): the two finite
half-translation laws follow from commutation, and its value at
the vacuum is unchanged because \(Af_0=0\).
The universal \(\mathrm{AllSome}\) clause gives
\[
                 f=T_0\Phi=(T_0+AT_0)\Phi=f+Af.
 \tag{JIO.41}
\]
Thus \(Af=0\), and evaluation at \(w=0\) gives \(\ell(f)=0\).
Every vector outside a subspace is separated by an algebraic
linear functional, as follows by adjoining that vector to a
basis of the subspace and extending to a basis.  Therefore
\(f\in\Theta\).

The original \(\mathrm{IsThPVecOf}\) record supplies precisely
this initial theta vacuum and this \(\mathrm{IsThetaAt}\)
condition.  Its place model acts only in finite directions.
The original vacuum argument in TAE9, now with the theta property
just proved, gives its Fock-vacuum equations at every real place:
extend a real annihilation functional algebraically to all
curves, perturb the same place model as in (JIO.41), and use
commutation with the finite Heisenberg action and annihilation
of \(f_0\).

CW12 proves that real Heisenberg derivatives of theta functions
are again theta functions, with genuine Banach derivatives.
The original scaled creators are fixed complex linear
combinations of two such derivatives.  They therefore preserve
\(\Theta\).  Their mutual commutation follows from the
Heisenberg commutator and the positive Cartan complex structure:
the two scalar commutators in
\([D_z-iD_{Jz},D_w-iD_{Jw}]\) cancel, including the cross terms.
Equivalently this is TAE3's Heisenberg calculus.
The finite monomial definition of \(\mathrm{pureTensorO}\)
therefore remains in \(\Theta\), proving every assertion of
(JIO.39).  In particular this proves the whole
\(R.\mathrm{thP}_{\mathrm{pt\_theta}}\) predicate; it is not
retained as a new input to the simultaneous construction.
% END JIO13

% BEGIN JIO14
\paragraph{JIO14: the original splitting characters have the same exponents for every \(t\).}
\label{claim:JIO14}
Let \(L\) be the original level-character data and assume
\(L.\mathrm{Is}(\mu_0)\).  For every \(i,t,\tau\), the characters
in the original pair
\(\mathrm{chrTh}(L.\mathrm{mu}(t),L.\eta,i+1)\) satisfy
\[
 \mathrm{HasInfExp}(\chi_V,\tau,-e_{i,\tau}),\qquad
 \mathrm{HasInfExp}(\chi_W,\tau,-e_{i,\tau}-\tfrac12).
 \tag{JIO.42}
\]
This is the exact \texttt{InfTypesAt} input needed below,
with no restriction on \(t\).

\emph{Proof.}
We first justify the actual character-model transport, rather
than identifying infinity exponents merely by name.
Suppose \(\mathrm{ChrEq}(c,\chi)\) and \(\chi\) has the
specified integral phase exponents \(n_\tau=2a_\tau\).
The definition of \(\mathrm{ChrEq}\) says that \(c\) is continuous
and trivial on principal ideles, and that its value on the
single finite-place principal element \(k\in E^\times\)
equals the original \(\chi.\mathrm{locK}_w(k)\).

These comparisons determine the real component as follows.
The finite tensor algebra
\(\mathbb A_{F,f}\otimes_F E\) has its usual components
\(\prod_{w\mid v}E_w\) over a place \(v\) of \(F\).
Indeed write \(E=F[\alpha]\); factoring its separable minimal
polynomial over \(F_v\) and applying the Chinese remainder
theorem gives this decomposition.  Outside the finite set
dividing its index and discriminant the corresponding integral
lattices agree, so the decompositions respect the restricted
product topology.  They give the original one-place idempotents
and the real/finite decomposition used in \(\mathrm{ChrEq}\).
The real tensor algebra is similarly
\(\prod_{\tau\in I}\mathbb C\), with the specified CM-type
embeddings; HCU14 proves this original coordinate realization.

For a fixed \(k\in E^\times\), multiply its one-place finite
principal elements over finite sets of places.  This net
converges to the purely finite component of the principal
idele: any restricted-product neighborhood specifies only
finitely many components, and \(k,k^{-1}\) are integral outside
a finite set.  The same net in the ordinary adelic model of
\(\chi\) has the same local values by \(\mathrm{ChrEq}\).
Continuity equates the two finite-component values.
Both full characters are trivial on the principal idele, so
their purely real values on \(k\) agree.  The latter is
\[
           \prod_\tau
              \left(\frac{\tau(k)}{|\tau(k)|}\right)^{n_\tau}
 \tag{JIO.43}
\]
by the original \(\mathrm{HasInfQ}\) condition.

The image of \(E\) in \(\prod_\tau\mathbb C\) is dense:
combine an \(F/\mathbb Q\) basis with the \(E/F\) basis
\(1,\delta\), where \(\overline\delta=-\delta\ne0\).
Their images form a real basis by the CM real tensor
isomorphism, and rational coordinates are dense in real
coordinates.  Restricting to the open subset
where all components are nonzero gives density of \(E^\times\).
Thus continuity extends (JIO.43) to every real idele.
Every infinity family \(\theta\) satisfying
\(\mathrm{IsInfFam}\) gives these same coordinates, since it
agrees on \(E\) and is continuous; equivalently apply the
same density argument to its component maps.
Taking a real idele supported at \(\tau\) proves the literal
\(\mathrm{HasInfExp}(c,\tau,a_\tau)\), including its universal
infinity-family clause and its existence clause.
The required coordinate family itself is supplied by the
real tensor decomposition just described.

Apply this to \(L.\mathrm{mu0}\) and \(L.\eta\) through the
actual fields \(h_L.\mathrm{mu0}\) and \(h_L.\eta\).
Their stored character-frame readings give exponents
\(-e_{i,\tau}\) and \(1/2\), respectively; reindexing the
CM type uses the specified bijection \(D.eTau\).
For \(L.\xi(t)\), the original field \(\xi_{\rm fin}\)
says that its Hecke model has finite order.  The real ideles
form the connected group \(\prod_\tau\mathbb C^\times\).
A continuous character of finite order has image in a finite
root-of-unity group, so is \(1\) there.  Hence its infinity
exponent is \(0\), and the same transport proves this for
the actual \(L.\xi(t)\).

The definition \(L.\mathrm{mu}(t)=\mathrm{tup}
(L.\mathrm{mu0},L.\xi(t))\) multiplies these two characters.
Multiplying and inverting their phase formulas adds and
negates the exponents; the integer phase exponents are
respectively added and negated as well.
Finally the actual pair is
\(\chi_V=L.\mathrm{mu}(t)_{i+1}\) and
\(\chi_W=L.\mathrm{mu}(t)_{i+1}L.\eta^{-1}\).
This gives (JIO.42) for every \(t\).  The finite order used
here may depend on \(t\); its only consequence is the exact
zero real exponent, so no parameter-dependent weight is
introduced.
% END JIO14

% BEGIN JIO15
\paragraph{JIO15: real determinant roots and their character values.}
\label{claim:JIO15}
For any original real-supported unitary \(g\) on a nondegenerate
Hermitian space \(S\), there is an idele \(x\), with finite
component \(1\), satisfying the actual determinant-root equation
\[
                  x=\det(\mathrm{mt}(g))\,\overline{x}.
 \tag{JIO.44}
\]
For every character \(\nu\) with integral infinity exponents
\(c_\tau\in\mathbb Z\), this same root can be chosen so that
\[
                         \nu(x)=\prod_\tau\det(g_\tau)^{c_\tau}.
 \tag{JIO.45}
\]

\emph{Proof.}
The relation \(g_\tau^*H_\tau g_\tau=H_\tau\), and
\(\det H_\tau\ne0\), imply \(|\det g_\tau|=1\).
Choose \(z_\tau\in\mathbb C\) with
\(z_\tau^2=\det g_\tau\) and \(|z_\tau|=1\), for example by
halving an argument.  The real tensor chart of HCU14 constructs
the idele whose real components are \(z_\tau\) and whose finite
component is \(1\); the same construction with \(z_\tau^{-1}\)
is its inverse.  Since \(z_\tau/\overline{z_\tau}=z_\tau^2\),
the root equation holds in every component and hence in the
actual tensor algebra.

Decompose this idele into the finite product of its one-place
real factors.  By definition \(\mathrm{HasInfExp}(\nu,\tau,c_\tau)\)
evaluates that factor as \(z_\tau^{2c_\tau}\).
The factor of two is essential: the source records twice the
rational infinity exponent as its integral phase exponent.
Multiplication gives
\(\prod_\tau z_\tau^{2c_\tau}
 =\prod_\tau(\det g_\tau)^{c_\tau}\).
Replacing a chosen square root by its negative does not change
this value, since \(2c_\tau\) is even.  This proves both the
original root contract and its simultaneous evaluation property.
% END JIO15

% BEGIN JIO16
\paragraph{JIO16: the exact scalar of the actual Weil vacuum.}
\label{claim:JIO16}
Let \(S,T,\psi,\theta,\omega,\chi_S,\chi_T\) be the original
objects satisfying \(\mathrm{IsWeilDatumO}\) and
\(\mathrm{IsInfFam}\), with the original
\(S.\mathrm{SigOK}\) and \(T.\mathrm{SigOK}\).
Suppose \(S\) has Cartan involution \(1\),
\(s_\tau\) is a Cartan involution on \(T\) with positive and
negative dimensions \(p_\tau,q_\tau\), and
\[
 \begin{gathered}
 \mathrm{HasInfExp}(\chi_S,\tau,a_\tau),\quad
 \mathrm{HasInfExp}(\chi_T,\tau,b_\tau),\\
 x_{1,\tau}=a_\tau+n/2,\qquad
 x_{2,\tau}=a_\tau-n/2,\qquad
 x_{3,\tau}=-(b_\tau+(p_\tau-q_\tau)/2)
                       \quad\text{are integers}.
 \end{gathered}
 \tag{JIO.46}
\]
For every original theta function \(f\) which is a Fock vacuum
for \(J_\tau z=i s_\tau z\) at every real place, and every
real-supported \(g,h\) with \(h_\tau s_\tau=s_\tau h_\tau\),
\[
 \omega(g,h)f=
 \left(\prod_\tau A_\tau^{x_{1,\tau}}
                     B_\tau^{x_{2,\tau}}
                     G_\tau^{x_{3,\tau}}\right)f ,
 \quad
 \begin{cases}
 A_\tau=\det(P_{+,\tau}h_\tau+P_{-,\tau}),\\
 B_\tau=\det(P_{-,\tau}h_\tau+P_{+,\tau}),\\
 G_\tau=\det g_\tau.
 \end{cases}
 \tag{JIO.47}
\]
This is the actual \texttt{vacuum\_character} scalar, not a
scalar inferred from covariance alone.

\emph{Proof.}
Use the explicit auxiliary character \(\eta_0\) and action
\(\omega_0\) constructed in CW24 on these same Hermitian
spaces and theta functions.  Its real exponent is \(1/2\);
its restriction to the base-field ideles is the particular
quadratic character of HF9, and HF11 constructs its Hecke
extension.  CW24 proves
exactly the operator construction needed here without assuming
that this auxiliary character satisfies the universal
\(\mathrm{IsEtaTilde}\) restriction.
The two supplied \(\mathrm{ResEta}\) fields can be tested
against that particular quadratic character.  Consequently
\[
          \nu_S=\chi_S\eta_0^{-n},\qquad
          \nu_T=\chi_T\eta_0^{-m}
 \tag{JIO.48}
\]
are trivial on all base-field ideles.  CW24's determinant-root
construction therefore gives a full Weil datum
\[
       \omega'(g,h)=\widetilde{\nu_T}(g)^{-1}
                          \widetilde{\nu_S}(h)\omega_0(g,h)
 \tag{JIO.49}
\]
with the prescribed original \(\chi_S,\chi_T\).
Here the tildes are evaluated on any roots in (JIO.44);
root independence, continuity, rational triviality and the
crossed finite splitting factor are proved in CW24 through
its explicit determinant-root argument.
BTH14 then identifies \(\omega'\) and the supplied \(\omega\)
on the whole original theta space.  This step uses two full
data with the same characters, and so fixes the actual action
normalization.

For a real-supported pair, the finite part of \(\omega_0\)
is the identity.  CW11--CW13 compute its scalar on the entire
vacuum space as the positive-Cartan determinant for the Hom
matrix \(h_\tau\otimes(g_\tau^{-1})^{\mathsf t}\).
Since \(S\) is positive, the positive subspace of that Hom
space is \(T_{\tau,+}\otimes S_\tau^*\).
The determinant on this subspace, with the identity on its
complement, is
\[
                          A_\tau^{\,n}G_\tau^{-p_\tau}.
 \tag{JIO.50}
\]
One can check this without choosing a new frame: conjugate
the idempotent \(P_+\) to \(\operatorname{diag}(1_{p},0_q)\).
The nonidentity block and its determinant are
\[
 h_+\otimes(g^{-1})^{\mathsf t},
 \qquad \det(h_+)^n\det(g^{-1})^p.
\]
The same block decomposition gives \(\det h=A_\tau B_\tau\).
All these determinants are nonzero.

The infinity exponents of the two characters in (JIO.48) are
\[
           a_\tau-n/2=x_{2,\tau},\qquad
           b_\tau-m/2=-x_{3,\tau}-p_\tau ,
 \tag{JIO.51}
\]
which are integers.  Evaluate the determinant roots by JIO15.
Multiplying their values with (JIO.50) gives at every place
\[
 G_\tau^{\,x_{3,\tau}+p_\tau}
       (A_\tau B_\tau)^{x_{2,\tau}}
       A_\tau^nG_\tau^{-p_\tau}
       =A_\tau^{x_{1,\tau}}B_\tau^{x_{2,\tau}}G_\tau^{x_{3,\tau}},
 \tag{JIO.52}
\]
since \(x_1=x_2+n\).
This proves (JIO.47).  The inverse factor evaluates
\(\nu_T\) on the source determinant root, while the direct
factor evaluates \(\nu_S\) on the target determinant root,
exactly as in (JIO.49).
The argument applies simultaneously to every vacuum \(f\);
it does not select a scalar separately for each vector.
No identification of \(\eta_0\) with the stored level
character \(L.\eta\) has been used.
% END JIO16

% BEGIN JIO17
\paragraph{JIO17: the original orthonormal-frame substitution.}
\label{claim:JIO17}
Let \(B_S^*H_SB_S=1\), let
\(B_T=\mathrm{onb}(M)=M.\mathrm{fr}^{-1}\), and let
\(s=\mathrm{cart}(M)=B_TJ_\epsilon B_T^{-1}\).
For \(k_\tau\in K_{\epsilon_\tau}\), \(u_\tau\in U(n)\)
and every theta function \(f\),
\[
 \mathrm{pt}_{B_Su,B_Tk}(f,P)
              =\mathrm{pt}_{B_S,B_T}(f,f_{\rm sub}(k,u)P).
 \tag{JIO.53}
\]
All creators have the original scale
\(2\sqrt{2\pi|\tau(\xi_0)|}\).

\emph{Proof.}
Write \(D_z\) for the real Heisenberg derivative and
\[
       C_J(z)=
        [2\sqrt{2\pi|\tau(\xi_0)|}]^{-1}(D_z-iD_{Jz}).
 \tag{JIO.54}
\]
The chosen host element \(\xi_0\ne0\) has nonzero image under
every embedding, so this scale is a fixed positive number.
The derivative is real-linear in \(z\), by CW12 or the explicit
Heisenberg calculus of TAE3.
If \(Jz=i\epsilon z\), direct expansion of
\(c=\Re c+i\Im c\) gives
\[
             C_J(cz)=
                  \begin{cases}c\,C_J(z),&\epsilon=+,\\
                                \overline c\,C_J(z),&\epsilon=-.
                   \end{cases}
 \tag{JIO.55}
\]
Indeed \(J(iz)=iJz\); substituting into (JIO.54) leaves \(c\)
in the positive case and \(\overline c\) in the negative one.

The exact original unit direction is
\(\mathrm{unitMap}(B_S,B_T;i,j)=B_TE_{ij}B_S^{-1}\).
Matrix multiplication gives
\[
 \mathrm{unitMap}(B_Su,B_Tk;i,j)
       =\sum_{a,b} k_{ai}(u^{-1})_{jb}
                      \mathrm{unitMap}(B_S,B_T;a,b).
 \tag{JIO.56}
\]
The latter direction has Cartan eigenvalue \(i\epsilon_a\).
If the coefficient is nonzero, sign preservation of \(k\)
gives \(\epsilon_a=\epsilon_i\).
For the negative case, unitarity converts
\(\overline{k_{ai}(u^{-1})_{jb}}\) to
\((k^{-1})_{ia}u_{bj}\).
Thus (JIO.55)--(JIO.56) produce exactly the two coefficients
of (JIO.5), with the same scale on both sides.

By JIO13 the creators preserve theta functions and commute
there.  Their polynomial evaluation is therefore independent
of a monomial's chosen ordering and satisfies
\(\mathrm{pt}(X_vP)=C_J(v)\mathrm{pt}(P)\).
Induct on a polynomial: the constant case is scalar
multiplication of \(f\), addition is linear, and the case
\(X_vP\) uses (JIO.56) with (JIO.55) and the induction
hypothesis.  This proves (JIO.53) on the actual pointwise
functions, without any completion or comparison scalar.
% END JIO17

% BEGIN JIO18
\paragraph{JIO18: the full actual compact Weil law for every polynomial.}
\label{claim:JIO18}
Under the original hypotheses of JIO16, use the orthonormal
source frames and model frames of JIO17.  Suppose \(g,h\)
are real-supported and
\[
       k_\tau=M.\mathrm{fr}_\tau\,h_\tau M.\mathrm{fr}_\tau^{-1}
                         \in K_{\epsilon_\tau},\qquad
       u_\tau=B_{S,\tau}^{-1}g_\tau B_{S,\tau}.
 \tag{JIO.57}
\]
Then every \(u_\tau\) is unitary and
\[
             \omega(g,h)\mathrm{pt}(f,P)
                           =\mathrm{pt}(f,F(k,u)P)
                   \quad\text{for every }P\in\mathcal P.
 \tag{JIO.58}
\]
This is the complete original
\texttt{om\_pureTensorO\_eq\_fact} contract.

\emph{Proof.}
The identity \(B_S^*H_SB_S=1\) makes \(B_S\) invertible
and gives \(H_S=(B_S^{-1})^*B_S^{-1}\).
Substitute this into \(g^*H_Sg=H_S\) to get \(u^*u=1\).
It also proves that \(1\) is a positive Cartan on \(S\).
The condition on \(k\) says that \(h\) commutes with
\(s=B_TJ_\epsilon B_T^{-1}\).  Its trace is
\(p-q\), so JIO16 applies to the same Cartan as the creators.

We check differentiation through the actual Weil operator.
Its Heisenberg covariance is
\[
              \omega(g,h)H_z=H_{hzg^{-1}}\omega(g,h).
\]
On theta functions, CW12 proves convergence of the Heisenberg
difference quotients in the Banach norm and hence in the
original compact-quotient \(L^2\) norm.  The original
\(\mathrm{om\_unit}\) makes \(\omega(g,h)\) an isometry for
that norm.  Passing to the difference-quotient limit proves
\[
             \omega(g,h)D_zf
                            =D_{hzg^{-1}}\omega(g,h)f .
 \tag{JIO.59}
\]
Both sides are theta functions.  Equality in \(L^2\) gives
pointwise equality: a nonzero continuous difference has
nonzero absolute value on an open set of positive Haar
measure on the compact additive quotient.
This argument does not assume continuity of the source's
ambient algebraic endomorphism on all functions.

Because \(h\) commutes with \(s\), the linear map
\(z\mapsto hzg^{-1}\) commutes with the Hom Cartan complex
structure.  Thus (JIO.59) intertwines the creators with
source frame \(gB_S=B_Su\) and target frame \(hB_T=B_Tk\).
Induction on monomials, each time using theta preservation,
gives
\[
 \omega(g,h)\mathrm{pt}_{B_S,B_T}(f,P)
              =\mathrm{pt}_{B_Su,B_Tk}(\omega(g,h)f,P).
 \tag{JIO.60}
\]
Use JIO17, then JIO16.  Conjugation by the stored frames
identifies the two target determinants in (JIO.47) with
\(d_+(k_\tau)\) and \(d_-(k_\tau)\); the source determinant
is \(\det u_\tau=\det g_\tau\).
The right side of (JIO.60) is consequently
\[
          c(k,u)\,\mathrm{pt}_{B_S,B_T}
                              (f,f_{\rm sub}(k,u)P)
                         =\mathrm{pt}_{B_S,B_T}(f,F(k,u)P),
 \tag{JIO.61}
\]
using linearity in the initial function and in \(P\).
This proves the exact formula, including all determinant
twists and every polynomial degree.
% END JIO18

% BEGIN JIO19
\paragraph{JIO19: the full simultaneous producer.}
\label{claim:JIO19}
Fix one original
\(R:\mathrm{RealDat}(D,A,K_s)\), on the selected Bundle if
one has been selected, and assume exactly
\[
 R.\mathrm{toLevCh}.\mathrm{Is}(A.\mu_0),\quad
 R.\mathrm{IsFam},\quad R.\mathrm{IsModels},\quad
 R.\mathrm{thP}_{\omega P},\quad R.\mathrm{thP}_{vP}.
 \tag{JIO.62}
\]
Then \(R.\mathrm{thP}_{\mathrm{pt\_theta}}\) holds, and
\[
 \begin{split}
 \forall i\ \exists\,\mathrm{io}_i\ \forall t:\quad
 \mathrm{IsIoW}\bigl(&R.\theta,\ R.MW(i,\mathrm{sg0}),\
        R.cV(i),\ R.\omega P(i,t),\\
                   &R.\mathrm{ptP}(i,t),\
                     \mathrm{RealDat}.\Delta(i),\
                     \mathrm{io}_i\bigr).
 \end{split}
 \tag{JIO.63}
\]
Equivalently a replacement field \(\mathrm{io}\) satisfies
\((\{R\text{ with ioP}:=\mathrm{io}\}).\mathrm{thP}_{\mathrm{ioP}}\).
This is the full original \texttt{JointHarm.exists\_ioP}.

\emph{Proof.}
JIO13 already proves the first conclusion.
Fix \(i\), but do not choose \(t\).
CFB13 installs on this same \(R,i\) the exact original
\(M=R.MW(i,\mathrm{sg0})\), \(c=R.cV(i)\), \(c_0=c_{0L,i}\),
\(B_S=R.BV(i)\), the literal coefficient weights and the
original \(p,q,a^+,a^-\).  Its arithmetic gives
\(m\le n\), \(p+q=m\), the ordered signs and both
antitonicity conditions.  Its integer twists satisfy both
identities (JIO.31) and all three equalities
\[
 x_{1,\tau}=-e_{i,\tau}+n/2,\qquad
 x_{2,\tau}=-e_{i,\tau}-n/2,\qquad
 x_{3,\tau}=-\bigl((p_\tau-q_\tau)/2-e_{i,\tau}-1/2\bigr).
 \tag{JIO.64}
\]
This installation uses the original \(\mathrm{IsFam.tie.mapL}\)
and \(\mathrm{IsModels.BV}\), and the exact base identity
\(\mathrm{sg0}=\mathrm{embLL}\).  No HolShape claim at another
coefficient embedding is involved.
Apply JIO12 once to obtain \(\mathrm{io}_i\), with the exact
unit-vector normalization.  It depends on \(i\) and the
stored data; no \(t\) has yet appeared.

Now let \(t\) be arbitrary.  JIO14 gives the two exponents
\(a=-e_i\), \(b=-e_i-1/2\) for the actual characters of
\(R.\omega P(i,t)\).  Equations (JIO.64) are exactly the
three integer equalities of JIO16.  JIO13 gives the theta
vacuum \(f=R.vP(i,t)\).  Hence JIO18 applies to the actual
\(\mathrm{ptP}(i,t)\).

For the first \(\mathrm{IsIoW}\) law, take an arbitrary
\(g\in G_\infty(V)\) and \(h=1\).
Then \(k_\tau=1\), and JIO18 makes every \(u_\tau\) unitary.
Use (JIO.36) in (JIO.58) to obtain
\[
 \omega(g,1)\mathrm{ptP}(i,t)(\mathrm{io}_i(x,w))
       =\mathrm{ptP}(i,t)(\mathrm{io}_i(x,c.\rho(g)w)).
 \tag{JIO.65}
\]
For the second law, fix \(\tau\) and \(k\in K_{\epsilon_\tau}\).
The matrix \(M.\mathrm{fr}_\tau^{-1}kM.\mathrm{fr}_\tau\)
preserves the original target Gram form.  HCU14's one-place
real lift is exactly the original
\[
 h=\mathrm{atTau}\bigl(\theta,W,\tau,
                M.\mathrm{fr}_\tau^{-1}kM.\mathrm{fr}_\tau\bigr).
 \tag{JIO.66}
\]
It is real-supported; its \(\tau\)-component is the displayed
matrix and all other components are \(1\).
In (JIO.57), therefore, the entire \(k\)-tuple is
\(\operatorname{single}_\tau k\).  Taking \(g=1\) gives
the entire \(u\)-tuple \(1\).
JIO18 and (JIO.37) yield the original target compact
intertwining identity on the updated lowest tensor.
This verifies the frame inverses and both component
identities explicitly.

The third \(\mathrm{IsIoW}\) law is (JIO.38); it contains
no \(t\), so its same witnesses work for every \(t\).
We have proved \(\exists\mathrm{io}_i\,\forall t\), in that
order.  Choosing these maps for all \(i\) gives the
replacement field in the statement.

Finally, if the original \(R\) has the weak stored-choice
law
\[
 \forall i\ \mathrm{OrZero}\bigl(
       \mathrm{PIoP}(i),R.\mathrm{ioP}(i)\bigr),\qquad
            \mathrm{PIoP}(i,y)\equiv\forall t\
                            \mathrm{IsIoW}(\ldots,t,\ldots,y),
 \tag{JIO.67}
\]
then (JIO.63) activates its first implication and proves
the strong predicate for the \emph{stored} \(R.\mathrm{ioP}(i)\).
The rule is existence followed by the given weak-choice
implication; it asserts no equality between the constructed
map and the stored map.
% END JIO19

\paragraph{Application boundary.}
Thus the simultaneous \(\mathrm{ioP}\) supplier is closed on
the exact original premises (JIO.62).  Combined with HCU20
and the inherited TAE13 packet, this removes that supplier
from the analytic extension interface for the same \(R\).
It does not supply the remaining bounded extremal form,
the actual target finite-volume pair, or the pointwise
smooth/Sobolev realization after Hilbert completion.
It also does not construct the selected Bundle, the original
coefficient tie, readings or rational tests; their own
producer contracts remain separate.
